\section{Transport of remainders and finite smoothness}
\label{sec:transport}

The theory so far concerns an exact finite germ. In practice the forward function is often known only through a truncated expansion, and one wants to know how the omitted part affects the inverse; conversely, one wants to turn a small residual $f(G(y))-y$ of an approximation $G$ into a bound for $G-g$. Both are mean-value arguments, but the scaling factor $z^{1-p}$ of the derivative must not be forgotten.

\subsection{From the forward function to the inverse}

\begin{theorem}[Transport of a forward remainder]
\label{thm:transport}
Let $f$ be the germ \eqref{eq:model} with $p>0$ and let $F=f+R$ on $(0,u_0)$, where
\begin{equation}
 R(u)=\OO\bigl(u^\rho\M(u)^k\bigr),\qquad R'(u)=\OO\bigl(u^{\rho-1}\M(u)^k\bigr),\qquad \rho>p,\ k\in\N .
 \label{eq:R-hypothesis}
\end{equation}
Then $F$ has a local inverse $G$ on the same branch as $g$, and
\begin{equation}
 G(y)-g(y)=\OO\bigl(z^{\rho-p+1}\M(z)^k\bigr)=\OO\bigl(w^{(\rho-p+1)/p}\M(w)^k\bigr),\qquad w=|y-y_0| .
 \label{eq:transport}
\end{equation}
The same conclusion holds if the derivative hypothesis is replaced by the assumption that $F$ is monotone near $0$.
\end{theorem}
\begin{proof}
From \eqref{eq:R-hypothesis} and Lemma~\ref{lem:dominance}, $F(u)-y_0\sim au^p$ and $F'(u)\sim apu^{p-1}$, so $F$ is strictly monotone near $0$ and its inverse $G$ satisfies $G\sim z$ (Theorem~\ref{thm:branch} applies verbatim). Put $u=G(y)$, $v=g(y)$; both are $\sim z$. Then $F(v)-y=F(v)-f(v)=R(v)=\OO(z^\rho\M(z)^k)$, while $F(u)-y=0$. On the interval between $u$ and $v$, which lies in $[z/2,2z]$ for small $z$, $F'$ is $\ge\frac12|ap|z^{p-1}2^{-|p-1|}$ in absolute value, so the mean value theorem gives $|u-v|\le|R(v)|/\min|F'|=\OO(z^{\rho-p+1}\M(z)^k)$. Under monotonicity instead of the derivative bound, use $f$'s derivative: $f(u)-f(v)=y-R(u)-y=-R(u)$, and $f'\asymp z^{p-1}$ on $[z/2,2z]$ gives the same estimate.
\end{proof}

\begin{remark}
The derivative hypothesis is not a consequence of the value bound: $R(u)=u^2\sin(u^{-2})$ is $\OO(u^2)$, but $R'(u)=2u\sin(u^{-2})-2u^{-1}\cos(u^{-2})$ destroys the monotonicity of $u+R(u)$, and $u+R$ has no inverse near $0$. Truncated expansions produced by \wl{Series} of log-exp functions do satisfy \eqref{eq:R-hypothesis}, because termwise differentiation of their asymptotic expansions is valid; the package records the hypothesis as \wl{"InputRemainder" -> \{$\rho$, k\}} and never proves it.
\end{remark}

\begin{corollary}[Cutoff cap]
If the inverse of $F$ is computed from the truncated germ $f$ with an exponent cutoff $q$ in $w$, the result is a valid expansion of $G$ only for $q\le(\rho-p+1)/p$; for such $q$ the remainder class is the coarser of the truncation remainder (Corollary~\ref{cor:remainder-y}) and \eqref{eq:transport}, compared by exponent first and by logarithmic degree at equal exponents.
\end{corollary}

For example $\sin x+x^2\log x=x+x^2\log x-x^3/6+\OO(x^5)$ is a germ with $p=1$, gaps $1,2$ and $\rho=5$, $k=0$; the inverse can be requested up to the cutoff $q=5$ and is $y-y^2\log y+y^3(\tfrac16+\log y+2\log^2y)+\OO(y^4\M^3)$ for $q=4$.

\subsection{From a residual to the error of the inverse}

\begin{theorem}[Residual bracket]
\label{thm:bracket}
Let $F$ be continuous on $[A,B]$ and differentiable inside with $F'\ge\mu>0$. Let $x_0\in(A,B)$, $y\in\R$, $\varrho=F(x_0)-y$. If $[x_0-|\varrho|/\mu,\ x_0+|\varrho|/\mu]\subseteq[A,B]$, then $F(x)=y$ has exactly one solution $x_*$ in $[A,B]$, it lies in that interval, and
\begin{equation}
 \frac{|\varrho|}{\max F'}\ \le\ |x_*-x_0|\ \le\ \frac{|\varrho|}{\mu} .
 \label{eq:bracket}
\end{equation}
\end{theorem}
\begin{proof}
$F(x_0+|\varrho|/\mu)\ge F(x_0)+|\varrho|\ge y$ and $F(x_0-|\varrho|/\mu)\le y$; the intermediate value theorem gives a root in the interval, monotonicity its uniqueness in $[A,B]$, and the mean value theorem both inequalities.
\end{proof}

This gives an a posteriori certificate that needs no expansion theory, only a lower bound on the derivative. For the two examples the global bounds \eqref{eq:global-slopes} give, for every $y>0$ and every $G>0$,
\begin{equation}
 |G-g_1(y)|\le\frac{|G+G^2(1+\log G)-y|}{1-2e^{-5/2}},\qquad |G-g_2(y)|\le|G+G^{\sqrt2}-y| .
 \label{eq:certificates}
\end{equation}
For a general germ the derivative is $\asymp z^{p-1}$, and the bracket must be applied on $[z/2,2z]$ with $\mu=\frac12|ap|z^{p-1}2^{-|p-1|}$ (valid for small $z$), giving $|G-g|\le\frac{2^{1+|p-1|}}{|ap|}z^{1-p}|f(G)-y|$.

\paragraph{The residual of the truncated expansion.} Combining Theorem~\ref{thm:remainder} with the mean value theorem, the truncated inverse $G_H$ satisfies
\begin{equation}
 \frac{f(G_H(y))}{a\,z^p}-1=-p\,z^{\sigma_*}C_{\sigma_*}(L)+\oo(z^{\sigma_*}),
 \label{eq:residual-identity}
\end{equation}
since $f(G_H)-y\approx f'(g)(G_H-g)=-apz^{p-1}\cdot z^{1+\sigma_*}C_{\sigma_*}$. The normalized residual has the same first block as the error, up to the factor $-p$; the package's \wl{InverseResidual} computes it exactly in the jet algebra (with a higher cutoff it displays the first omitted block; below the cutoff it must vanish identically, which is an independent check of every coefficient), and its ``forward residual cutoff'' in the variable of $y$ is $q+(p-1)/p$.

\subsection{Finite smoothness suffices}

The finite germ is analytic away from $0$, and the convergence proof used that. For perturbations outside the finite power--log class (a coefficient $\sin(\log x)$, a term $e^{-1/x}$, a function known only through derivative bounds) the marker formula \eqref{eq:near-identity} can still be applied, and the following theorem, which appears in report~05 with a sketched proof, shows that it is a genuine asymptotic expansion whenever the perturbation and finitely many of its derivatives are small in the power--log scale. We give the complete proof.

\begin{theorem}[Tame near-identity inversion]
\label{thm:tame}
Let $N\ge0$, $\delta>0$ and $M\ge0$. Let $\varphi$ be real-valued and $C^{N+1}$ on $(0,x_0)$ with
\begin{equation}
 \varphi^{(j)}(x)=\OO\bigl(x^{1+\delta-j}\M(x)^M\bigr)\qquad(x\to0^+),\quad j=0,1,\dots,N+1 .
 \label{eq:tame-hypothesis}
\end{equation}
Then $f=x+\varphi$ has a unique local inverse $g$ on $(0,y_0)$ with $g(y)\sim y$, and
\begin{equation}
 g(y)=y+\sum_{n=1}^N\frac{(-1)^n}{n!}\frac{\dd^{n-1}}{\dd y^{n-1}}\bigl[\varphi(y)^n\bigr]+\OO\bigl(y^{1+(N+1)\delta}\M(y)^{(N+1)M}\bigr).
 \label{eq:tame}
\end{equation}
\end{theorem}
\begin{proof}
\emph{Step 1: the homotopy.} For $\lambda\in[0,1]$ put $F_\lambda(x)=x+\lambda\varphi(x)$. By \eqref{eq:tame-hypothesis} with $j=1$, $\varphi'(x)\to0$, so there is $x_1$ with $F_\lambda'\ge\frac12$ on $(0,x_1)$ for all $\lambda\in[0,1]$; each $F_\lambda$ is an increasing bijection of $(0,x_1)$ onto an interval containing $(0,y_1)$ for some $y_1$ independent of $\lambda$. Let $G(\lambda,y)$ be its inverse, so that
\begin{equation}
 G+\lambda\varphi(G)=y,\qquad G(0,y)=y .
 \label{eq:homotopy}
\end{equation}
Since $|\varphi(x)|\le x/2$ for $x<x_1$ (after shrinking $x_1$), $F_\lambda(y/2)<y<F_\lambda(2y)$ and therefore $G(\lambda,y)\in[y/2,2y]$ for all $\lambda$ and all small $y$. The map $(\lambda,x)\mapsto x+\lambda\varphi(x)-y$ is $C^{N+1}$ with $x$-derivative $\ge\frac12$, so by the implicit function theorem $G$ is $C^{N+1}$ in $(\lambda,y)$, with
\begin{equation}
 \partial_yG=\frac1{1+\lambda\varphi'(G)}\in\bigl[\tfrac23,2\bigr],\qquad \partial_\lambda G=-\frac{\varphi(G)}{1+\lambda\varphi'(G)}=\psi(G)\,\partial_yG,\quad\psi:=-\varphi .
 \label{eq:transport-pde}
\end{equation}

\emph{Step 2: the transport identity.} For every $C^{n}$ function $\Psi$ and $1\le n\le N+1$,
\begin{equation}
 \partial_\lambda^n\bigl[\Psi(G)\bigr]=\partial_y^{n-1}\bigl[\psi(G)^n\,\partial_y\Psi(G)\bigr].
 \label{eq:transport-identity}
\end{equation}
For $n=1$: $\partial_\lambda\Psi(G)=\Psi'(G)\partial_\lambda G=\Psi'(G)\psi(G)\partial_yG=\psi(G)\partial_y[\Psi(G)]$. Assume \eqref{eq:transport-identity} for $n$ and all $\Psi$. Then, using it once more for the functions $\psi^n$ and $\Psi$ in the second step,
\begin{align*}
 \partial_\lambda^{n+1}[\Psi(G)]&=\partial_y^{n-1}\,\partial_\lambda\bigl[\psi(G)^n\partial_y\Psi(G)\bigr]
 =\partial_y^{n-1}\bigl[\psi(G)\,\partial_y(\psi(G)^n)\,\partial_y\Psi(G)+\psi(G)^n\,\partial_y\bigl(\psi(G)\partial_y\Psi(G)\bigr)\bigr]\\
 &=\partial_y^{n-1}\,\partial_y\bigl[\psi(G)^{n+1}\partial_y\Psi(G)\bigr],
\end{align*}
where the last equality is the product rule read backwards: $\partial_y[\psi^{n+1}\Psi_y]=(n+1)\psi^n\psi_y\Psi_y+\psi^{n+1}\Psi_{yy}$ and the bracket equals $n\psi^n\psi_y\Psi_y+\psi^n(\psi_y\Psi_y+\psi\Psi_{yy})$, the same thing. (Here $\psi_y$ stands for $\partial_y[\psi(G)]$.)

\emph{Step 3: the coefficients.} At $\lambda=0$, $G=y$, so \eqref{eq:transport-identity} with $\Psi=\mathrm{id}$ gives $\partial_\lambda^nG(0,y)=\partial_y^{n-1}[\psi(y)^n]=(-1)^n\partial_y^{n-1}[\varphi(y)^n]$, the Lagrange terms. Taylor's formula in $\lambda$ with integral remainder gives
\[
 g(y)=G(1,y)=y+\sum_{n=1}^N\frac{(-1)^n}{n!}\partial_y^{n-1}[\varphi^n]+\frac1{N!}\int_0^1(1-\lambda)^N\,\partial_\lambda^{N+1}G(\lambda,y)\,\dd\lambda .
\]

\emph{Step 4: the remainder.} Put $\kappa(y)=y^\delta\M(y)^M$. Since $G\in[y/2,2y]$ and $\M(G)\le2\M(y)$, hypothesis \eqref{eq:tame-hypothesis} gives $\varphi^{(j)}(G)=\OO(\kappa\,y^{1-j})$ uniformly in $\lambda$, for $j\le N+1$. We claim that for $1\le j\le N+1$,
\begin{equation}
 \partial_y^jG=\OO(y^{1-j})\quad(j\ge1),\qquad \partial_y^j[\varphi(G)]=\OO(\kappa\,y^{1-j})\quad(j\ge0),
 \label{eq:derivative-bounds}
\end{equation}
uniformly in $\lambda\in[0,1]$. Differentiating \eqref{eq:homotopy} $j$ times in $y$ and solving for $\partial_y^jG$: $(1+\lambda\varphi'(G))\partial_y^jG$ equals $\delta_{j1}$ minus $\lambda$ times a sum of terms $\varphi^{(m)}(G)\prod\partial_y^{i_s}G$ with $2\le m$, $\sum i_s=j$, $i_s\ge1$ (Fa\`a di Bruno). By induction on $j$ each such term is $\OO(\kappa y^{1-m})\prod\OO(y^{1-i_s})=\OO(\kappa y^{1-j})$, so $\partial_y^jG=\OO(\kappa y^{1-j})\subseteq\OO(y^{1-j})$ for $j\ge2$ (and $\partial_yG=\OO(1)$). The same Fa\`a di Bruno expansion with $m\ge1$ gives the second claim. By Leibniz, $\partial_y^j[\psi(G)^{N+1}]=\OO(\kappa^{N+1}y^{N+1-j})$, and $\partial_y^{N-j}[\partial_yG]=\partial_y^{N-j+1}G=\OO(y^{-(N-j)})$. Applying \eqref{eq:transport-identity} with $\Psi=\mathrm{id}$ and $n=N+1$,
\[
 \partial_\lambda^{N+1}G=\partial_y^N\bigl[\psi(G)^{N+1}\partial_yG\bigr]=\sum_{j=0}^N\binom Nj\,\partial_y^j[\psi(G)^{N+1}]\,\partial_y^{N-j+1}G=\OO\bigl(\kappa^{N+1}y\bigr)
\]
uniformly in $\lambda$, and the integral remainder is $\OO(y\kappa(y)^{N+1})=\OO(y^{1+(N+1)\delta}\M(y)^{(N+1)M})$.
\end{proof}

\begin{remark}
Applied to $\varphi(x)=x^2(1+\log x)$ ($\delta=1$, $M=1$) the theorem re-proves the remainder $\OO(y^{N+2}\M^{N+1})$ of \eqref{eq:log-remainder} with a proof that uses no analyticity at all; applied to $\varphi(x)=x^2\sin(\log x)$ ($\delta=1$, $M=0$) it justifies $g(y)=y-y^2\sin L+y^3\sin L(2\sin L+\cos L)+\OO(y^4)$, $L=\log y$, a perturbation whose coefficients oscillate and which no power--log expansion describes. The package's \wl{PerturbativeInverse} produces the terms; the theorem is the user's tool for asserting the remainder.
\end{remark}
